\sheet{4}{MPI processes and message protocols}{1h30}
\section{Preparation}
Define \textbf{process}, \textbf{rank}, \textbf{communicator}, \textbf{tag},
\textbf{collective}, \textbf{request} and \textbf{deadlock}. Draw two processes
and their separate arrays. Explain why changing one array does not change the
other automatically, and why four ranks do not imply four machines.

Sources are in \code{04-mpi}. D1--D4 are independent programs;
D5 measures D2. Every file has its own MPI initialization, numerical checks and
\code{main}. Use the MPI C interface from C and matching compiler/launcher
installations. Commands run from this \code{04-mpi} directory:
\begin{lstlisting}[style=shell]
cmake -S . -B build -DCMAKE_BUILD_TYPE=Release
\end{lstlisting}
On a cluster, use its allocation and launcher settings. Rebuild after edits.
If a run hangs, stop it and compare the message/collective sequence across ranks
before trying a larger input.

Each C source is a self-contained program: arrays, allocation, checks and
timing are visible in that file. No course-specific header is needed. Complete
the functions marked TODO; an unfinished function reports its name and exits.
The numerical arguments are optional; their ranges and defaults are listed in
\code{README.md}. Each run checks the requested case. Run small and edge cases
explicitly by changing the arguments, then measure representative larger inputs.

\section{Exercises}
\exercise{D1}{Distributing an uneven input}
Edit \code{D1-distribute.c}. Rank 0 initially owns the complete input.
For global length $N$ and $P$ ranks, use counts and displacements
\[n_r=\lfloor N/P\rfloor+[r<N\bmod P],\qquad d_0=0,\quad d_{r+1}=d_r+n_r.\]
The supplied \code{partition} function fills the count and displacement arrays.
The caller's local length and start are \code{counts[rank]} and \code{displs[rank]}. The bracket contributes one when its condition is true.
\begin{enumerate}
\item Draw and tabulate counts, displacements and owned global indices for
$N=10,P=3$ and $N=2,P=4$. Verify coverage without overlap. Explain the empty blocks.
\item Implement \code{distribute}. The supplied main allocates the arrays and
passes the local length, counts and displacements. Call \code{MPI_Scatterv}. Root supplies the global elements, counts and displacements;
every rank supplies its own receive storage and count.
\item Pass the array pointer as the buffer argument: send its elements, not the
bytes of the pointer variable itself. For a double datatype, displacement 4 selects element
4, not byte 4. Every rank participates even when its count is zero.
\item Build and run:
\begin{lstlisting}[style=shell]
cmake --build build --target d1_distribute -j
mpiexec -n 3 ./build/bin/d1_distribute 10
mpiexec -n 4 ./build/bin/d1_distribute 2
\end{lstlisting}
The argument is global length. The program checks local values and prints each
rank's count and start. Also run empty global input and uneven blocks.
Explain why root-only preparation does not imply a root-only collective call.
\end{enumerate}

\exercise{D2}{Local calculation, reduction and gathering}
Edit \code{D2-vectors.c}. This separate program distributes both $x$ and $y$,
computes $z_i=2x_i+1$ locally and returns the global dot product on every rank.
It gathers $z$ at rank 0. Reuse your distribution reasoning from D1 in this file.
\begin{enumerate}
\item Implement \code{distributed_vectors}. Obtain the partition, allocate
local storage and scatter $x$ and then $y$. Calculate a local affine vector and
partial dot product using ordinary loops. Empty blocks contribute sum zero.
\item Combine one scalar per rank with \code{MPI_Allreduce} and addition.
Return the resulting scalar from \code{distributed_vectors} on every rank.
\item The supplied main allocates \code{gathered} only at root. Use
\code{MPI_Gatherv} with the original counts and displacements to fill this
output array. Other ranks receive a \code{NULL} output argument, ignored by MPI.
\item Write the common sequence: scatter $x$, scatter $y$, all-reduce, gather.
Explain why all ranks, including empty ones, must follow it compatibly.
\item Build and run:
\begin{lstlisting}[style=shell]
cmake --build build --target d2_vectors -j
mpiexec -n 3 ./build/bin/d2_vectors 17 3
\end{lstlisting}
The arguments are global length and trials. Checks verify the dot product on
every rank and every gathered element. Repeat with 1, 2 and 4 ranks. For
$x=(0,1,2,3,4)$ and three ranks, predict the counts and gathered $z$ before running.
\end{enumerate}

\exercise{D3}{A blocking ring exchange}
Edit \code{D3-ring-blocking.c}. The C function receives two output pointers,
\code{from_left} and \code{from_right}; write their pointed-to values. Rank $r$ sends one value to both neighbours
and obtains their original values. In a ring, left is $(r+P-1)\bmod P$ and
right is $(r+1)\bmod P$.
\begin{enumerate}
\item Draw three ranks with initial values $r+1$. Label one set of arrows
sending right and another sending left. Use a distinct tag for each direction.
Predict rank 0's two received values.
\item Implement \code{ring_exchange} with two \code{MPI_Sendrecv} calls.
Send right while receiving from the left, then reverse directions. Keep the
original send value separate from both received outputs.
\item Build and run:
\begin{lstlisting}[style=shell]
cmake --build build --target d3_ring -j
mpiexec -n 3 ./build/bin/d3_ring
\end{lstlisting}
The program takes no numerical arguments and checks several input values.
Repeat with one and two ranks: explain matching when a neighbour is the caller
itself or when both neighbours are the same process.
\item Explain why a standard blocking send on every rank followed by receives
can deadlock. Success with a small buffered message does not establish a correct
protocol for larger messages.
\end{enumerate}

\exercise{D4}{A nonblocking ring exchange}
Edit \code{D4-ring-nonblocking.c}. This is a separate executable with the same
input/output contract as D3. A nonblocking start returns before completion is
established; buffer access rules continue to apply after that return.
\begin{enumerate}
\item Implement \code{ring_exchange} by posting two receives and two sends.
Use the same neighbours and tag convention as D3. Retain all four requests.
\item Call \code{MPI_Waitall} before returning. Keep send storage alive and
unchanged until both sends complete. Do not access either receive buffer until
its receive completes. Both sends may read the same unchanged original value.
\item Build and run:
\begin{lstlisting}[style=shell]
cmake --build build --target d4_ring -j
mpiexec -n 3 ./build/bin/d4_ring
\end{lstlisting}
Repeat with one, two and four ranks. A completed send makes its local buffer
reusable; it does not establish that the remote application has finished using
the received data.
\item Draw the interval from initiation to completion for each buffer. Identify
what unrelated computation could be done safely in that interval. This ring has
little useful work to overlap; a nonblocking name alone does not imply speedup.
\end{enumerate}

\exercise{D5}{Measuring the distributed vector calculation}
Use \code{d2_vectors}; no additional source file is needed. The timed region
includes local allocation, both scatters, local arithmetic, all-reduction and
gathering. Input generation, checking, printing and timing aggregation are excluded.
Each trial reports the maximum rank interval; the summary takes the median across
those maxima. One initial run is discarded as warm-up.
\begin{enumerate}
\item Keep the global length fixed and run with 1, 2 and 4 ranks within the
allocation, for example:
\begin{lstlisting}[style=shell]
mpiexec -n 4 ./build/bin/d2_vectors 1048576 7
\end{lstlisting}
Record whether the ranks run on one or several machines, their placement, median
and range. Use the same build and numerical requirements.
\item Calculate $S_P=T_1/T_P$ and $E_P=S_P/P$ relative to the same MPI program
with one rank. Plot time or speedup. Explain which communication costs remain
when local arithmetic shrinks.
\item Repeat at a second fixed global length. Explain why a timing of only
rank 0's arithmetic would answer a different question. Do not infer network cost
or multi-node performance from a localhost result alone.
\end{enumerate}
